\pdfminorversion=4
\documentclass[aspectratio=169]{beamer}

\mode<presentation>
{
  \usetheme{default}
  \usecolortheme{default}
  \usefonttheme{default}
  \setbeamertemplate{navigation symbols}{}
  \setbeamertemplate{caption}[numbered]
  \setbeamertemplate{footline}[frame number]  % or "page number"
  \setbeamercolor{frametitle}{fg=white}
  \setbeamercolor{footline}{fg=black}
} 

\usepackage[english]{babel}
\usepackage[utf8x]{inputenc}
\usepackage{tikz}
\usepackage{courier}
\usepackage{array}
\usepackage{bold-extra}
\usepackage{minted}
\usepackage[thicklines]{cancel}
\usepackage{fancyvrb}

\xdefinecolor{dianablue}{rgb}{0.18,0.24,0.31}
\xdefinecolor{darkblue}{rgb}{0.1,0.1,0.7}
\xdefinecolor{darkgreen}{rgb}{0,0.5,0}
\xdefinecolor{darkgrey}{rgb}{0.35,0.35,0.35}
\xdefinecolor{darkorange}{rgb}{0.8,0.5,0}
\xdefinecolor{darkred}{rgb}{0.7,0,0}
\definecolor{darkgreen}{rgb}{0,0.6,0}
\definecolor{mauve}{rgb}{0.58,0,0.82}

\title[2019-05-22-pattern-matching-irishep]{Pattern matching for particle physics}
\author{Jim Pivarski}
\institute{Princeton University -- IRIS-HEP}
\date{May 22, 2019}

\usetikzlibrary{shapes.callouts}

\begin{document}

\logo{\pgfputat{\pgfxy(0.11, 7.4)}{\pgfbox[right,base]{\tikz{\filldraw[fill=dianablue, draw=none] (0 cm, 0 cm) rectangle (50 cm, 1 cm);}\mbox{\hspace{-8 cm}\includegraphics[height=1 cm]{princeton-logo-long.png}\hspace{0.1 cm}\raisebox{0.1 cm}{\includegraphics[height=0.8 cm]{iris-hep-logo-long.png}}\hspace{0.1 cm}}}}}

\begin{frame}
  \titlepage
\end{frame}

\logo{\pgfputat{\pgfxy(0.11, 7.4)}{\pgfbox[right,base]{\tikz{\filldraw[fill=dianablue, draw=none] (0 cm, 0 cm) rectangle (50 cm, 1 cm);}\mbox{\hspace{-8 cm}\includegraphics[height=1 cm]{princeton-logo.png}\hspace{0.1 cm}\raisebox{0.1 cm}{\includegraphics[height=0.8 cm]{iris-hep-logo.png}}\hspace{0.1 cm}}}}}

% Uncomment these lines for an automatically generated outline.
%\begin{frame}{Outline}
%  \tableofcontents
%\end{frame}

% START START START START START START START START START START START START START

\begin{frame}{}
\large
\vspace{1.25 cm}
Quite a few groups have been thinking about physics event processing languages, explicitly or implicitly.

\vspace{0.25 cm}
\begin{itemize}
\item \textcolor{darkblue}{LHADA/ADL:} Sezen Sekmen, Harry Prosper, Philippe Gras
\item \textcolor{darkblue}{CutLang:} Gokhan Unel
\item \textcolor{darkblue}{IRIS-HEP Analysis Systems:} Gordon Watts, Mason Proffitt, Emma Torro
\item \textcolor{darkblue}{FAST-Carpenter (YAML):} Benjamin Krikler
\item \textcolor{darkblue}{NAIL:} Andrea Rizzi
\item \textcolor{darkblue}{RDataFrame:} Enrico Guiraud, Danilo Piparo, and the ROOT Team
\item \textcolor{darkblue}{AEACuS \& RHADAManTHUS:} Joel Walker (phenomenology)
\item \textcolor{darkblue}{Femtocode:} me, though not for several years\ldots
\end{itemize}

\vspace{0.25 cm}
\small \hfill see \textcolor{blue}{\href{https://indico.cern.ch/event/769263/timetable/}{Analysis Description Languages Workshop}} (May 6--8)
\end{frame}

\begin{frame}{}
\Large
\vspace{1.25 cm}
\begin{columns}
\column{0.87\linewidth}
\begin{center}
For a physics domain-specific language (DSL) to be useful, it has to solve \textcolor{darkblue}{physics problems} in a way that is clearly simpler than the general-purpose language.

\vspace{1 cm}
\uncover<2->{This is different from just being ``high-level,'' solving any problem with less baggage than a low-level language.}
\end{center}
\end{columns}
\end{frame}

\begin{frame}{What is difficult in physics analyses?}
\large
\vspace{0.5 cm}
\begin{itemize}\setlength{\itemsep}{0.35 cm}
\item<1-> Distributing work to remote processors. General enough for a whole industry. (Spark, Parsl, Dask, Condor, Slurm\ldots)
\item<1-> Interactive analysis, publication-quality plotting. Also very general. (Jupyter)
\item<1-> Workflow management: chaining analysis tasks. General. (CWL, Makefile)
\item<2-> Booking, filling, managing histograms. \textcolor{darkblue}{Physics specific!} (ROOT)
\item<2-> Complex, multi-component distribution fitting. \textcolor{darkblue}{Physics specific!} (RooFit)
\item<2-> Complex multi-histogram fitting. \textcolor{darkblue}{Physics specific!} (Combine, HistFactory)
\item<3-> Signal/control regions, systematic variations. \textcolor{darkblue}{Physics specific!} (slide 2)
\item<3-> Particle combinatorics: complex decay chains. \textcolor{darkblue}{Physics specific!} (???)
\end{itemize}
\end{frame}

\begin{frame}[fragile]{Examples of pattern matching}
\Large
\vspace{0.25 cm}
\textcolor{darkorange}{\bf Python:} (can only unpack iterables into tuples)

\begin{center}
\begin{minipage}{0.85\linewidth}
\normalsize
\begin{minted}{python}
def tree():
    return [[[1, 2], [3, 4]], [5, 6]]
(((a, b), c), _) = tree()    # a → 1, b → 2, c → [3, 4]
\end{minted}
\end{minipage}
\end{center}

\vspace{0.35 cm}
\begin{uncoverenv}<2->
\textcolor{darkorange}{\bf Haskell:}

\begin{center}
\begin{minipage}{0.85\linewidth}
\normalsize
\begin{minted}{haskell}
pz :: (Particle particle) => particle -> Float
pz (Neutral pt eta _)   = pt * sinh(eta)
pz (Charged pt eta _ q) = pt * sinh(eta)
\end{minted}
\end{minipage}
\end{center}
\end{uncoverenv}

\vspace{0.35 cm}
\begin{uncoverenv}<3->
\textcolor{darkorange}{\bf Scala:}

\begin{center}
\begin{minipage}{0.85\linewidth}
\normalsize
\begin{minted}{scala}
def pz(particle: Particle) = match particle {
    case Neutral(pt, eta, _)    => pt * sinh(eta)
    case Charged(pt, eta, _, q) => pt * sinh(eta)
}
\end{minted}
\end{minipage}
\end{center}
\end{uncoverenv}
\end{frame}

\end{document}
